\documentclass[11pt,a4paper]{article}
\usepackage[utf8]{inputenc}
\usepackage[T1]{fontenc}
\usepackage[french]{babel}
\usepackage{amsmath,amssymb}
\usepackage{geometry}
\usepackage{enumitem}
\usepackage{booktabs}
\usepackage{array}
\usepackage{tabularx}
\usepackage[table]{xcolor}
\usepackage{tcolorbox}
\usepackage{fancyhdr}
\usepackage{multicol}
\usepackage{tikz}
\usepackage{pgfplots}
\pgfplotsset{compat=1.18}
\usepackage{lastpage}

\tcbuselibrary{skins,breakable}

\geometry{margin=1.9cm, headheight=15pt}

% Couleurs
\definecolor{primary}{RGB}{31,78,121}
\definecolor{soft}{RGB}{230,239,247}
\definecolor{amber}{RGB}{178,91,0}
\definecolor{amberbg}{RGB}{255,244,224}
\definecolor{amberborder}{RGB}{240,192,112}
\definecolor{greenbg}{RGB}{233,246,234}
\definecolor{greenborder}{RGB}{182,224,184}
\definecolor{redbg}{RGB}{253,236,236}
\definecolor{redborder}{RGB}{242,179,171}
\definecolor{green}{RGB}{46,125,50}
\definecolor{red}{RGB}{211,47,47}
\definecolor{grayline}{RGB}{120,130,140}

% En-tête et pied de page
\pagestyle{fancy}
\fancyhf{}
\fancyhead[L]{\textcolor{primary}{\textbf{BTS BioALC — 2\textsuperscript{ème} année}}}
\fancyhead[R]{\textcolor{primary}{Fiche de travail — Contrôle qualité et synthèse}}
\fancyfoot[L]{\textcolor{grayline}{\small Nom : \dotfill}}
\fancyfoot[C]{\textcolor{grayline}{\small \thepage/\pageref{LastPage}}}
\fancyfoot[R]{\textcolor{grayline}{\small Semaine du 16/11/2026}}
\renewcommand{\headrulewidth}{0.4pt}
\renewcommand{\footrulewidth}{0.4pt}

% Boîtes
\newtcolorbox{consigne}{
  colback=soft, colframe=primary, boxrule=0.8pt, arc=4pt,
  title={\textbf{Mode d'emploi}}, fonttitle=\bfseries,
  coltitle=white, colbacktitle=primary,
  left=6pt, right=6pt, top=4pt, bottom=4pt,
  breakable
}

\newtcolorbox{contexte}{
  colback=amberbg, colframe=amberborder, boxrule=0.8pt, arc=4pt,
  title={\textbf{Contexte professionnel}}, fonttitle=\bfseries,
  coltitle=black, colbacktitle=amberborder,
  left=6pt, right=6pt, top=4pt, bottom=4pt,
  breakable
}

\newtcolorbox{rappel}[1][]{
  colback=greenbg, colframe=greenborder, boxrule=0.8pt, arc=4pt,
  fonttitle=\bfseries, coltitle=black, colbacktitle=greenborder,
  left=6pt, right=6pt, top=4pt, bottom=4pt,
  breakable, #1
}

\newtcolorbox{questionbox}[1][]{
  colback=white, colframe=primary, boxrule=0.8pt, arc=4pt,
  fonttitle=\bfseries, coltitle=white, colbacktitle=primary,
  left=6pt, right=6pt, top=4pt, bottom=4pt,
  breakable, #1
}

\newtcolorbox{sommative}[1][]{
  colback=redbg, colframe=redborder, boxrule=1pt, arc=4pt,
  fonttitle=\bfseries, coltitle=white, colbacktitle=red,
  left=8pt, right=8pt, top=6pt, bottom=6pt,
  breakable, #1
}

\newtcolorbox{sommativeTT}[1][]{
  colback=soft, colframe=primary, boxrule=1pt, arc=4pt,
  fonttitle=\bfseries, coltitle=white, colbacktitle=primary,
  left=8pt, right=8pt, top=6pt, bottom=6pt,
  breakable, #1
}

\newtcolorbox{autoeval}{
  colback=redbg, colframe=redborder, boxrule=0.8pt, arc=4pt,
  title={\textbf{Auto-évaluation}}, fonttitle=\bfseries,
  coltitle=black, colbacktitle=redborder,
  left=6pt, right=6pt, top=4pt, bottom=4pt,
  breakable
}

% Commandes utilitaires
\newcommand{\lignes}[1]{\vspace{0.15cm}\foreach \x in {1,...,#1}{\noindent\rule{\linewidth}{0.4pt}\par\vspace{0.35cm}}}
\newcommand{\case}{\raisebox{0.1ex}{\fbox{\rule{0pt}{1.1ex}\rule{1.1ex}{0pt}}}}

\title{\textcolor{primary}{\textbf{Fiche de travail — Contrôle qualité et synthèse statistique}}\\
\large Semaine du 16 au 20 novembre 2026 — BTS Bioanalyses en Laboratoire de Contrôle}
\author{}
\date{}

\begin{document}

\maketitle

\begin{consigne}
Cette fiche couvre la semaine 3 du module « Statistique descriptive ». Elle a deux objectifs :
\begin{itemize}[leftmargin=*]
    \item \textbf{Relier} les outils statistiques aux \textbf{enseignements professionnels} (contrôle qualité, validation de méthodes).
    \item \textbf{Synthétiser} l'ensemble des indicateurs vus depuis le début du module.
\end{itemize}

La dernière partie contient une \textbf{évaluation sommative notée} proposée en deux versions :
\begin{itemize}[leftmargin=*]
    \item \textbf{Version A — standard} : durée \textbf{20 minutes} ;
    \item \textbf{Version B — tiers temps} : durée \textbf{27 minutes} (soit + 1/3 du temps standard).
\end{itemize}
La somme des deux durées (20 + 27 = \textbf{47 minutes}) reste \textbf{inférieure à 50 minutes}, ce qui permet de faire passer les deux versions dans le même créneau horaire sans dépassement.
\end{consigne}

\vspace{0.3cm}

\begin{contexte}
Le responsable qualité du laboratoire vous confie une mission : \textbf{valider la méthode de dosage du glucose} utilisée en routine. Il vous fournit un jeu de données de contrôle et vous demande de produire un \textbf{rapport statistique} conforme aux exigences du contrôle qualité interne (CIQ).
\end{contexte}

\vspace{0.4cm}

% ============================================================
\section*{Partie 1 — Lien avec les enseignements professionnels}
% ============================================================

\subsection*{1.1. Le contrôle qualité en laboratoire}

\begin{rappel}[title={\textbf{À retenir — Le contrôle qualité interne (CIQ)}}]
Le contrôle qualité interne consiste à analyser régulièrement un \textbf{matériau de contrôle} (échantillon dont on connaît la valeur cible) et à vérifier que les résultats restent \textbf{dans les limites acceptables}. On utilise pour cela :
\begin{itemize}[leftmargin=*]
    \item la \textbf{moyenne} (valeur cible) ;
    \item l'\textbf{écart-type} (dispersion acceptée) ;
    \item les \textbf{limites} \(\pm 2\sigma\) (alerte) et \(\pm 3\sigma\) (rejet) ;
    \item la \textbf{carte de contrôle} pour suivre la stabilité dans le temps.
\end{itemize}
\end{rappel}

\begin{questionbox}[title={\textbf{Question 1 — Compréhension}}]
Expliquez en 3 à 4 lignes pourquoi un laboratoire doit effectuer un contrôle qualité interne chaque jour.

\lignes{4}
\end{questionbox}

\subsection*{1.2. Validation de méthodes}

\begin{rappel}[title={\textbf{À retenir — Validation de méthode}}]
Avant d'utiliser une méthode en routine, il faut \textbf{prouver} qu'elle est fiable. Les indicateurs statistiques utilisés sont :
\begin{itemize}[leftmargin=*]
    \item \textbf{justesse} : la moyenne est-elle proche de la valeur vraie ? (biais faible)
    \item \textbf{fidélité} : les résultats sont-ils reproductibles ? (écart-type faible)
    \item \textbf{répétabilité} : dispersion intra-laboratoire ;
    \item \textbf{reproductibilité} : dispersion inter-laboratoires.
\end{itemize}
\end{rappel}

\begin{questionbox}[title={\textbf{Question 2 — Application}}]
Un laboratoire teste une nouvelle méthode de dosage. Sur un même échantillon de contrôle, il obtient les 10 résultats suivants (en g/L) :

\begin{center}
\begin{tabular}{|c|c|c|c|c|c|c|c|c|c|}
\hline
0,89 & 0,91 & 0,88 & 0,90 & 0,92 & 0,89 & 0,90 & 0,91 & 0,89 & 0,91 \\
\hline
\end{tabular}
\end{center}

La valeur cible est de \textbf{0,90 g/L}. L'écart-type maximal toléré est de \textbf{0,02 g/L}.

\vspace{0.2cm}
\textbf{a)} Calculez la \textbf{moyenne} des 10 résultats.

\lignes{3}

\textbf{b)} Calculez l'\textbf{écart-type} (utilisez le tableur ou faites le calcul à la main).

\lignes{3}

\textbf{c)} La méthode est-elle \textbf{juste} ? Est-elle \textbf{fidèle} ? Justifiez.

\lignes{4}

\textbf{d)} Rédigez la conclusion du rapport destiné au responsable qualité (3 à 4 lignes).

\lignes{4}
\end{questionbox}

\subsection*{1.3. Cartes de contrôle (introduction)}

\begin{rappel}[title={\textbf{À retenir — Carte de contrôle}}]
Une carte de contrôle est un graphique qui suit l'évolution d'un indicateur dans le temps. Trois lignes y figurent :
\begin{itemize}[leftmargin=*]
    \item \textbf{ligne centrale (LC)} = moyenne de référence ;
    \item \textbf{limite supérieure (LSC)} = LC + \(3\sigma\) ;
    \item \textbf{limite inférieure (LIC)} = LC \(-\) \(3\sigma\).
\end{itemize}
Un point \textbf{hors des limites} signale une dérive à investiguer.
\end{rappel}

\begin{questionbox}[title={\textbf{Question 3 — Lecture de carte}}]
Voici une carte de contrôle simplifiée des résultats de glucose obtenus sur 12 jours :

\begin{center}
\begin{tikzpicture}[scale=0.75]
  \draw[->, thick] (0,0) -- (13,0) node[right] {\small Jours};
  \draw[->, thick] (0,0) -- (0,6) node[above] {\small g/L};

  \draw[red, thick, dashed] (0,4.5) -- (13,4.5) node[right, font=\small] {LSC = 0,95};
  \draw[primary, thick] (0,3) -- (13,3) node[right, font=\small] {LC = 0,90};
  \draw[red, thick, dashed] (0,1.5) -- (13,1.5) node[right, font=\small] {LIC = 0,85};

  \fill[greenbg] (0,1.5) rectangle (13,4.5);

  \foreach \x/\y in {1/2.8, 2/3.1, 3/2.9, 4/3.2, 5/3.0, 6/2.7, 7/3.3, 8/3.1, 9/2.9, 10/3.0, 11/4.7, 12/3.1} {
    \fill[primary] (\x,\y) circle (0.1);
  }

  \fill[red] (11,4.7) circle (0.15);
  \node[above, red, font=\small] at (11,4.7) {Hors limite !};
\end{tikzpicture}
\end{center}

\textbf{a)} Combien de points sont hors des limites ?

\lignes{1}

\textbf{b)} À quel jour le problème est-il détecté ?

\lignes{1}

\textbf{c)} Que doit faire le technicien ce jour-là ?

\lignes{3}
\end{questionbox}

\newpage

% ============================================================
\section*{Partie 2 — Synthèse statistique descriptive}
% ============================================================

\subsection*{2.1. Récapitulatif des indicateurs}

\begin{rappel}[title={\textbf{Fiche de synthèse — Tous les indicateurs}}]
\begin{tabularx}{\linewidth}{|p{2.6cm}|X|p{4.2cm}|}
\hline
\rowcolor{soft}
\textbf{Indicateur} & \textbf{Formule / Méthode} & \textbf{Quand l'utiliser ?} \\ \hline
\textbf{Moyenne} & \(\bar{x} = \dfrac{1}{n}\sum x_i\) & Distribution symétrique, sans valeur aberrante. \\ \hline
\textbf{Médiane} & Valeur centrale après tri & Distribution asymétrique ou valeurs aberrantes. \\ \hline
\textbf{Mode} & Valeur la plus fréquente & Variables qualitatives ou discrètes. \\ \hline
\textbf{Étendue} & \(E = x_{\max} - x_{\min}\) & Aperçu rapide, très sensible aux extrêmes. \\ \hline
\textbf{Variance} & \(V = \dfrac{1}{n}\sum (x_i - \bar{x})^2\) & Comparaison avec la moyenne. \\ \hline
\textbf{Écart-type} & \(\sigma = \sqrt{V}\) & Même unité que les données. \\ \hline
\textbf{\(Q_1\), \(Q_3\)} & Quartiles après tri & Mesure robuste de dispersion. \\ \hline
\textbf{IQR} & \(IQR = Q_3 - Q_1\) & Détection de valeurs aberrantes. \\ \hline
\end{tabularx}
\end{rappel}

\subsection*{2.2. Choix du couple d'indicateurs}

\begin{questionbox}[title={\textbf{Exercice — Associer situation et couple d'indicateurs}}]
Pour chaque situation, indiquez le couple adapté :
\begin{itemize}[leftmargin=*]
    \item[\textbf{(i)}] (moyenne, variance)
    \item[\textbf{(ii)}] (médiane, IQR)
    \item[\textbf{(iii)}] (médiane, mode)
\end{itemize}

\vspace{0.2cm}
\begin{center}
\renewcommand{\arraystretch}{1.5}
\begin{tabularx}{\linewidth}{|X|c|}
\hline
\rowcolor{soft}
\textbf{Situation} & \textbf{Couple choisi} \\ \hline
1. Mesures de pH homogènes, un seul opérateur. & \hspace{2cm} \\ \hline
2. Salaires de 200 techniciens avec quelques très hauts salaires. & \\ \hline
3. Groupes sanguins de 500 donneurs. & \\ \hline
4. Dosage de glucose sur 30 patients, sans valeur extrême. & \\ \hline
5. Temps de réponse d'un automate, avec une panne ponctuelle. & \\ \hline
\end{tabularx}
\end{center}
\end{questionbox}

\subsection*{2.3. Rédaction d'un rapport statistique}

\begin{questionbox}[title={\textbf{Exercice — Rapport de contrôle qualité}}]
On vous fournit les indicateurs suivants pour deux séries de mesures de glucose (n = 20) :

\vspace{0.2cm}
\begin{center}
\begin{tabular}{|l|c|c|}
\hline
\rowcolor{soft}
\textbf{Indicateur} & \textbf{Série A (routine)} & \textbf{Série B (nouvelle méthode)} \\ \hline
Moyenne \(\bar{x}\) & 0,895 & 0,905 \\ \hline
Médiane \(Q_2\) & 0,885 & 0,910 \\ \hline
Écart-type \(\sigma\) & 0,045 & 0,078 \\ \hline
\(Q_1\) & 0,860 & 0,850 \\ \hline
\(Q_3\) & 0,930 & 0,960 \\ \hline
IQR & 0,070 & 0,110 \\ \hline
\end{tabular}
\end{center}

\vspace{0.3cm}
Rédigez un rapport de 5 à 6 lignes à destination du responsable qualité. Vous devez :
\begin{itemize}[leftmargin=*]
    \item comparer les deux séries ;
    \item identifier celle qui est la plus homogène ;
    \item conclure sur la fiabilité de la nouvelle méthode ;
    \item proposer une action.
\end{itemize}

\lignes{6}
\end{questionbox}

\newpage

% ============================================================
\section*{Évaluation sommative 1 — Version A (standard)}
% ============================================================

\begin{sommative}[title={\textbf{Évaluation sommative 1 — Version A (20 minutes)}}]
\textbf{Durée : 20 minutes — Calculatrice autorisée — Note sur 20.}

\vspace{0.2cm}
\textbf{Toutes les réponses doivent être justifiées.}

\vspace{0.4cm}
\textbf{Exercice 1 — QCM (5 points, 1 pt par question).} Entourez la bonne réponse.

\vspace{0.3cm}
\textbf{1.} La moyenne est un indicateur :
\begin{itemize}[leftmargin=1.2cm, itemsep=1pt]
    \item[A.] de position \hspace{2cm} B. de dispersion \hspace{2cm} C. de forme
\end{itemize}

\textbf{2.} L'écart-type est :
\begin{itemize}[leftmargin=1.2cm, itemsep=1pt]
    \item[A.] la racine carrée de la variance \hspace{2cm} B. le carré de la variance \hspace{2cm} C. la moitié de la variance
\end{itemize}

\textbf{3.} Pour la série \(4, 5, 5, 6, 7, 7, 7, 8, 9, 10\), la médiane vaut :
\begin{itemize}[leftmargin=1.2cm, itemsep=1pt]
    \item[A.] 5 \hspace{2cm} B. 6,5 \hspace{2cm} C. 7
\end{itemize}

\textbf{4.} L'écart interquartile est :
\begin{itemize}[leftmargin=1.2cm, itemsep=1pt]
    \item[A.] \(Q_2 - Q_1\) \hspace{2cm} B. \(Q_3 - Q_1\) \hspace{2cm} C. \(Q_3 - Q_2\)
\end{itemize}

\textbf{5.} Une valeur est suspecte si elle sort de :
\begin{itemize}[leftmargin=1.2cm, itemsep=1pt]
    \item[A.] \([\bar{x} - \sigma\ ;\ \bar{x} + \sigma]\)
    \item[B.] \([Q_1 - 1{,}5 \times IQR\ ;\ Q_3 + 1{,}5 \times IQR]\)
    \item[C.] \([x_{\min}\ ;\ x_{\max}]\)
\end{itemize}

\vspace{0.4cm}
\textbf{Exercice 2 — Étude d'une série (8 points).}

On mesure le taux de cholestérol (en g/L) sur 10 patients :

\begin{center}
\begin{tabular}{|c|c|c|c|c|c|c|c|c|c|}
\hline
1,80 & 1,95 & 2,10 & 1,75 & 2,00 & 1,85 & 2,20 & 1,90 & 2,05 & 1,95 \\
\hline
\end{tabular}
\end{center}

\textbf{a)} Triez les valeurs par ordre croissant. \hfill \textbf{(1 pt)}

\lignes{2}

\textbf{b)} Calculez la moyenne \(\bar{x}\). \hfill \textbf{(2 pts)}

\lignes{3}

\textbf{c)} Déterminez la médiane \(Q_2\), puis \(Q_1\) et \(Q_3\). \hfill \textbf{(2 pts)}

\lignes{4}

\textbf{d)} Calculez l'écart interquartile \(IQR\). \hfill \textbf{(1 pt)}

\lignes{2}

\textbf{e)} Une nouvelle mesure donne 3,50 g/L. Est-elle aberrante ? Justifiez. \hfill \textbf{(2 pts)}

\lignes{3}

\vspace{0.4cm}
\textbf{Exercice 3 — Contrôle qualité (7 points).}

Un laboratoire utilise un matériau de contrôle dont la valeur cible est 5,00 g/L et l'écart-type toléré \(\sigma \leq 0{,}15\) g/L. Sur 8 jours, les résultats suivants ont été obtenus :

\begin{center}
\begin{tabular}{|c|c|c|c|c|c|c|c|c|}
\hline
Jour & 1 & 2 & 3 & 4 & 5 & 6 & 7 & 8 \\
\hline
Valeur (g/L) & 5,02 & 4,98 & 5,05 & 4,95 & 5,10 & 5,00 & 4,92 & 5,08 \\
\hline
\end{tabular}
\end{center}

\textbf{a)} Calculez la moyenne de ces 8 mesures. \hfill \textbf{(1 pt)}

\lignes{2}

\textbf{b)} Calculez l'écart-type. \hfill \textbf{(2 pts)}

\lignes{3}

\textbf{c)} La méthode est-elle \textbf{juste} ? \textbf{Fidèle} ? Justifiez avec les critères donnés. \hfill \textbf{(2 pts)}

\lignes{4}

\textbf{d)} Rédigez une conclusion de 3 lignes pour le responsable qualité. \hfill \textbf{(2 pts)}

\lignes{4}

\end{sommative}

\newpage

% ============================================================
\section*{Évaluation sommative 1 — Version B (tiers temps)}
% ============================================================

\begin{sommativeTT}[title={\textbf{Évaluation sommative 1 — Version B (tiers temps, 27 minutes)}}]
\textbf{Durée : 27 minutes — Calculatrice autorisée — Note sur 20.}

\vspace{0.2cm}
\textbf{Toutes les réponses doivent être justifiées. Les consignes sont explicitées pour vous aider à organiser votre travail.}

\vspace{0.4cm}
\textbf{Exercice 1 — QCM (5 points, 1 pt par question).} Entourez la bonne réponse.

\vspace{0.3cm}
\textbf{1.} La moyenne est un indicateur :
\begin{itemize}[leftmargin=1.2cm, itemsep=1pt]
    \item[A.] de position \hspace{2cm} B. de dispersion \hspace{2cm} C. de forme
\end{itemize}

\textbf{2.} L'écart-type est :
\begin{itemize}[leftmargin=1.2cm, itemsep=1pt]
    \item[A.] la racine carrée de la variance \hspace{2cm} B. le carré de la variance \hspace{2cm} C. la moitié de la variance
\end{itemize}

\textbf{3.} Pour la série \(3, 4, 4, 5, 6, 6, 6, 7, 8, 9\), la médiane vaut :
\begin{itemize}[leftmargin=1.2cm, itemsep=1pt]
    \item[A.] 5 \hspace{2cm} B. 6 \hspace{2cm} C. 6,5
\end{itemize}

\textbf{4.} L'écart interquartile est :
\begin{itemize}[leftmargin=1.2cm, itemsep=1pt]
    \item[A.] \(Q_2 - Q_1\) \hspace{2cm} B. \(Q_3 - Q_1\) \hspace{2cm} C. \(Q_3 - Q_2\)
\end{itemize}

\textbf{5.} Une valeur est suspecte si elle sort de :
\begin{itemize}[leftmargin=1.2cm, itemsep=1pt]
    \item[A.] \([\bar{x} - \sigma\ ;\ \bar{x} + \sigma]\)
    \item[B.] \([Q_1 - 1{,}5 \times IQR\ ;\ Q_3 + 1{,}5 \times IQR]\)
    \item[C.] \([x_{\min}\ ;\ x_{\max}]\)
\end{itemize}

\vspace{0.4cm}
\textbf{Exercice 2 — Étude d'une série (8 points).}

On mesure le taux de glucose (en g/L) sur 8 patients :

\begin{center}
\begin{tabular}{|c|c|c|c|c|c|c|c|}
\hline
0,90 & 0,95 & 1,00 & 0,85 & 1,05 & 0,92 & 0,88 & 0,95 \\
\hline
\end{tabular}
\end{center}

\textbf{a)} Triez les valeurs par ordre croissant. Écrivez-les ci-dessous en les séparant par des virgules. \hfill \textbf{(1 pt)}

\lignes{2}

\textbf{b)} Calculez la moyenne \(\bar{x}\). Rappel : \(\bar{x} = \dfrac{\text{somme des valeurs}}{\text{nombre de valeurs}}\). \hfill \textbf{(2 pts)}

\lignes{3}

\textbf{c)} Déterminez la médiane \(Q_2\), puis \(Q_1\) et \(Q_3\). Rappel : la médiane est la valeur du milieu après tri. \hfill \textbf{(2 pts)}

\lignes{4}

\textbf{d)} Calculez l'écart interquartile \(IQR = Q_3 - Q_1\). \hfill \textbf{(1 pt)}

\lignes{2}

\textbf{e)} Une nouvelle mesure donne 1,40 g/L. Est-elle aberrante ? Utilisez la règle \([Q_1 - 1{,}5 \times IQR\ ;\ Q_3 + 1{,}5 \times IQR]\). \hfill \textbf{(2 pts)}

\lignes{3}

\vspace{0.4cm}
\textbf{Exercice 3 — Contrôle qualité (7 points).}

Un laboratoire utilise un matériau de contrôle dont la valeur cible est 4,00 g/L et l'écart-type toléré \(\sigma \leq 0{,}10\) g/L. Sur 6 jours, les résultats suivants ont été obtenus :

\begin{center}
\begin{tabular}{|c|c|c|c|c|c|c|}
\hline
Jour & 1 & 2 & 3 & 4 & 5 & 6 \\
\hline
Valeur (g/L) & 4,02 & 3,98 & 4,05 & 3,96 & 4,01 & 3,99 \\
\hline
\end{tabular}
\end{center}

\textbf{a)} Calculez la moyenne de ces 6 mesures. \hfill \textbf{(1 pt)}

\lignes{2}

\textbf{b)} Calculez l'écart-type. Rappel : \(\sigma = \sqrt{\dfrac{1}{n}\displaystyle\sum (x_i - \bar{x})^2}\). \hfill \textbf{(2 pts)}

\lignes{4}

\textbf{c)} La méthode est-elle \textbf{juste} ? (la moyenne est-elle proche de la cible ?) \textbf{Fidèle} ? (l'écart-type respecte-t-il la tolérance ?) Justifiez. \hfill \textbf{(2 pts)}

\lignes{4}

\textbf{d)} Rédigez une conclusion de 3 lignes pour le responsable qualité. \hfill \textbf{(2 pts)}

\lignes{4}

\end{sommativeTT}

\newpage

% ============================================================
\section*{Bilan de la semaine}
% ============================================================

\begin{autoeval}
Cochez la case correspondant à votre niveau pour chaque compétence :

\vspace{0.3cm}
\begin{center}
\begin{tabularx}{\linewidth}{|X|c|c|c|}
\hline
\rowcolor{soft}
\textbf{Compétence} & \textbf{Acquis} & \textbf{En cours} & \textbf{À revoir} \\ \hline
Je sais expliquer le rôle du contrôle qualité interne. & \case & \case & \case \\ \hline
Je sais calculer moyenne et écart-type sur une série. & \case & \case & \case \\ \hline
Je sais déterminer \(Q_1\), \(Q_2\), \(Q_3\) et l'IQR. & \case & \case & \case \\ \hline
Je sais repérer une valeur aberrante. & \case & \case & \case \\ \hline
Je sais choisir le couple d'indicateurs adapté. & \case & \case & \case \\ \hline
Je sais lire une carte de contrôle simple. & \case & \case & \case \\ \hline
Je sais rédiger un rapport statistique court. & \case & \case & \case \\ \hline
\end{tabularx}
\end{center}

\vspace{0.3cm}
\textbf{Une question que je me pose encore après cette semaine :}

\lignes{2}
\end{autoeval}

\vspace{0.4cm}

\begin{contexte}
\textbf{Pour aller plus loin (facultatif, non noté) :} recherchez la norme \textbf{ISO 7870-2} (cartes de contrôle Shewhart) et citez en 5 lignes les règles de Nelson permettant de détecter une dérive.

\lignes{5}
\end{contexte}

\vspace{0.5cm}

\begin{center}
\textit{Fiche de travail — BTS BioALC — 2\textsuperscript{ème} année — Semaine du 16/11/2026}\\
\textit{Document à conserver pour les révisions et la préparation de l'examen.}
\end{center}

\end{document}